\section{竞争格局：技术领先、生态领先与国家级议题}

\subsection{从公司竞争到生态竞争}

在 Lex 提问 ``竞争优势'' 时，Jensen 并没有把答案停留在产品参数，而是回到 installed base、开发者速度和系统交付能力。这说明他把竞争对象定义为 ``生态系统''。当竞争单位从公司上升到生态，策略重心就会从短期发布节奏转向长期网络效应。

\begin{knowledgebox}{三层竞争结构}
可以把当前格局理解为三层：
\begin{itemize}
    \item 第一层：芯片和系统性能（产品层）。
    \item 第二层：开发者与软件栈（平台层）。
    \item 第三层：供应链与产业协同（基础设施层）。
\end{itemize}
这三层里，后两层更决定长期份额。
\end{knowledgebox}

\subsection{技术领导力与国家安全话题}

访谈后段出现 ``technology leadership'' 与 ``national security'' 的讨论。这里的重点不是口号，而是事实：算力平台已成为关键基础设施，影响科研、工业、金融、国防和公共服务。平台企业因此同时承受商业责任与公共责任。

\begin{warningbox}{叙事风险：把技术竞争简化为零和博弈}
技术领导力确实涉及国家利益，但如果把问题简化为单维对抗，会忽略全球供应链相互依赖这一现实。更可行的路径通常是：在安全边界内保持合作，在合作边界内保持竞争。
\end{warningbox}

\subsection{本章小结}

NVIDIA 的竞争方式已经从 ``产品竞争'' 延展到 ``生态和基础设施竞争''。这也解释了为什么管理层在公开叙事中同时谈技术、产业和政策。

